%% ma: L11-B03-0010
%% lop: 11
%% chuong: 1
%% bai: 3
%% dang: 11.3.7
%% loai: TN4
%% muc: NB
%% dap_an: "A"
%% nguon: "de-cuong-giua-ky-1-toan-11-nam-2026-2027-truong-thpt-hoang-van-thu-ha-noi-2-8.pdf (M08, Đề cương giữa kì 1 Toán 11 THPT Hoàng Văn Thụ, Hà Nội), Phần 1 Câu 12"
%% kien_thuc: "Bài 3 (hàm số lượng giác)"
%% trang_thai: cho-duyet
%% ly_do: "Dạng mới đề xuất 11.3.7"
\begin{ex}
Đường cong trong hình dưới đây là đồ thị của hàm số nào?
\begin{center}\begin{tikzpicture}[x=0.55cm,y=0.9cm]
\draw[truc] (-4.2,0)--(7.6,0) node[below] {$x$};
\draw[truc] (0,-1.5)--(0,1.6) node[left] {$y$};
\node[below left] at (0,0) {$O$};
\draw[phu] (-3.1416,0)--(-3.1416,-1)--(3.1416,-1)--(3.1416,0) (0,1)--(6.2832,1)--(6.2832,0);
\node[above] at (-3.1416,0) {$-\pi$}; \node[above] at (3.1416,0) {$\pi$}; \node[below] at (6.2832,0) {$2\pi$};
\node[above right] at (0,1) {$1$}; \node[below right] at (0,-1) {$-1$};
\draw[thick,smooth,samples=120,domain=-4.2:7.4] plot (\x,{cos(\x r)});
\end{tikzpicture}\end{center}
\choice{\True $y=\cos x$}{$y=1-\sin x$}{$y=1+\sin x$}{$y=\sin x$}
\loigiai{Đồ thị đi qua $(0;1)$, đạt giá trị nhỏ nhất $-1$ tại $x=\pm\pi$ và giá trị lớn nhất $1$ tại $x=0$, $x=2\pi$; đó là đồ thị hàm số $y=\cos x$.}
\end{ex}
